\documentclass[11pt]{article}
\usepackage[margin=1in]{geometry}
\usepackage{amsmath,amssymb,amsthm}
\usepackage[hidelinks]{hyperref}
\newtheorem{theorem}{Theorem}
\newtheorem{lemma}{Lemma}
\title{Semicompleteness of hypercube division sequences\\
\large A solution to MathDB 349741}
\author{Principia Math In Progress}
\date{18 September 2026}
\begin{document}
\maketitle
\begin{abstract}
For every positive dimension, the hypercube division sequence has a
stronger property than semicompleteness: every integer between zero and
the sum of any finite prefix is a sum of distinct terms of that prefix.
The proof applies the classical complete-sequence criterion to an exact
binomial identity. This proves Hoelscher's Conjecture 1.3 in the published
paper (Conjecture 1.0.3 in its arXiv version).
\end{abstract}

\section*{Task 1: the original problem and its full solution}
For integers $M,N\geq1$, let $C_M^N$ denote the maximum number of pieces
obtained by cutting an $M$-dimensional hypercube with $N-1$ hyperplanes.
For a fixed $M$, Hoelscher asks whether, for every $j>2$ and every positive
integer
\[
 k<\left\lfloor \frac{1}{j}\sum_{i=1}^j C_M^i\right\rfloor,
\]
one can express $k$ as a sum of distinct fractions chosen from
$C_M^1/j,\ldots,C_M^j/j$. The cases $M=1,2,3,4$ are already proved in
\cite[Theorem 1.2 and Section 3]{H}. We prove the assertion for all $M$.
The subset-sum criterion used below is classical, due to Brown
\cite{B}, and is also used in \cite[Section 3]{H}. We include its elementary
proof; neither that criterion nor the region-count formula is claimed
as new.

\begin{theorem}\label{thm:main}
Fix $M\geq1$ and $j\geq1$, and put $S_{M,j}=\sum_{i=1}^j C_M^i$.
For every integer $t$ with $0\leq t\leq S_{M,j}$ there is a subset
$I\subseteq\{1,\ldots,j\}$ such that
\[
 t=\sum_{i\in I}C_M^i.
\]
Consequently $\{C_M^N\}_{N\geq1}$ is semicomplete for every $M\geq1$.
\end{theorem}

\begin{lemma}[Finite-prefix form of the complete-sequence criterion]
Let $a_1,\ldots,a_j$ be positive integers, with $a_1=1$. Suppose
\[
 a_{n+1}\leq 1+\sum_{i=1}^n a_i\qquad(1\leq n<j).
\]
Then every integer from $0$ to $\sum_{i=1}^j a_i$ is the sum of a
subset of the indexed terms.
\end{lemma}
\begin{proof}
For one term the possible sums are $0,1$. Suppose the first $n$ terms
give every integer in $[0,S]$, where $S=\sum_{i=1}^n a_i$.
Subsets omitting or including $a_{n+1}$ give, respectively, all integers
in $[0,S]$ and $[a_{n+1},a_{n+1}+S]$. The hypothesis says that these
two integer intervals overlap or are adjacent. Their union is every
integer in $[0,S+a_{n+1}]$, proving the induction.
\end{proof}

\begin{proof}[Proof of Theorem \ref{thm:main}]
Use $\binom{n}{r}=0$ for $r>n$. The standard hyperplane-region formula is
\begin{equation}\label{eq:regions}
 C_M^N=\sum_{r=0}^M\binom{N-1}{r}.
\end{equation}
For completeness, the maximal number $R_M(n)$ of regions of $n$
hyperplanes in $\mathbb R^M$ satisfies
$R_M(n)=R_M(n-1)+R_{M-1}(n-1)$, with $R_M(0)=R_0(n)=1$:
the new hyperplane adds as many regions as the induced arrangement
on it. General position attains the upper bound. Pascal's identity
then gives $R_M(n)=\sum_{r=0}^M\binom nr$.
The same maximum is attained in a hypercube: choose an interior point
of each of the finitely many full-space regions and scale and translate
the entire arrangement so all these points lie inside the hypercube.
Convexity ensures that intersecting a region with the cube never
creates more than one piece. This proves \eqref{eq:regions}.

The hockey-stick identity gives, for every $N\geq1$,
\begin{equation}\label{eq:sum}
 S_{M,N}=\sum_{r=0}^M\sum_{i=1}^N\binom{i-1}{r}
        =\sum_{r=0}^M\binom N{r+1}.
\end{equation}
Subtracting \eqref{eq:regions}, with $N+1$ in place of $N$, yields
the exact identity
\begin{equation}\label{eq:gap}
 \boxed{\quad 1+S_{M,N}-C_M^{N+1}=\binom N{M+1}\geq0.\quad}
\end{equation}
Also $C_M^1=1$, so the lemma proves the asserted finite-prefix
representation. The terms are positive and strictly increasing, since
\[
 C_M^{N+1}-C_M^N=\sum_{r=0}^{M-1}\binom{N-1}{r}\geq1.
\]
Thus distinct indices give distinct numerators, as required in the source.

Finally, if $j>2$ and $0<k<\lfloor S_{M,j}/j\rfloor$, then
$0<kj<S_{M,j}$. Apply the proved representation to $t=kj$ and divide
by $j$. This gives precisely the required sum of distinct fractions.
In fact the representation holds for every integer
$0\leq k\leq\lfloor S_{M,j}/j\rfloor$.
\end{proof}

\paragraph{Constructive interpretation.}
Consider the terms in descending order, selecting a term exactly when
it does not exceed the current remainder. Suppose the remainder $t$
is at most $S_{M,n}$. At $a_n=C_M^n$, either subtract $a_n$ or leave
$t$ unchanged. The new remainder is at most $S_{M,n-1}$, because
\[
 t\geq a_n\ \Longrightarrow\ t-a_n\leq S_{M,n-1},
 \qquad
 t<a_n\ \Longrightarrow\ t\leq a_n-1\leq S_{M,n-1}.
\]
The final remainder is zero. No optimality in the number of selected
terms is asserted.

\section*{Tasks 2--5: proposed follow-up directions}
These are proposed questions, not additional solved claims.
\begin{description}
\item[Task 2.] For a fixed $M$, classify the indices $r$ for which deleting
the term $C_M^r$ and reindexing the remaining sequence preserves
semicompleteness.
\item[Task 3.] Let $a_N=\sum_{r=0}^d b_r\binom{N-1}{r}$, where
$b_r$ are nonnegative integers, $b_0\geq1$, and $b_1\geq1$.
Characterize the coefficient vectors for which $(a_N)$ is semicomplete,
and distinguish them from those for which every finite prefix has all
subset sums between zero and its total.
\item[Task 4.] For fixed $M$, obtain a uniform asymptotic formula in the
central range for the number of subsets $I\subseteq\{1,\ldots,j\}$
with $\sum_{i\in I}C_M^i=t$, as $j\to\infty$.
\item[Task 5.] For fixed $M$, determine the largest possible ratio,
over positive representable targets in a prefix, between the number
of terms selected by the descending algorithm above and the minimum
number of distinct terms needed for that target.
\end{description}

\paragraph{Verification and status.}
The accompanying exact verifier checks the region recurrence against
the closed form, enumerates finite subset sums, and reconstructs
representations. These finite checks supplement the proof for all
parameters. The published nine-page source and the current eight-page
arXiv version were read in full; both retain the general conjecture.
The dated literature audit reports bounded checks for earlier claims.
All verification is by the authoring agent; no independent peer review
is claimed.

\begin{thebibliography}{9}
\raggedright
\bibitem{H} Z. Hoelscher,
\emph{Semicomplete Arithmetic Sequences, Division of Hypercubes,
and the Pell Constant},
The PUMP Journal of Undergraduate Research \textbf{4} (2021), 108--116.
\href{https://doi.org/10.46787/pump.v4i0.2524}{doi:10.46787/pump.v4i0.2524};
\href{https://arxiv.org/abs/2102.07083}{arXiv:2102.07083v1}.
\bibitem{B} J. L. Brown, Jr.,
\emph{Note on Complete Sequences of Integers},
American Mathematical Monthly \textbf{68} (1961), 557--560.
\href{https://doi.org/10.2307/2311150}{doi:10.2307/2311150}.
\end{thebibliography}
\end{document}
